\begin{pastexam}[label={px:ccu-113-swa-10}]{中正資工所甲組 113 學年度 第 15 題}
（單選）

(2\%) The solution of the recurrence $\mathrm{T}(n)=3\mathrm{T}\left(\frac{n}{5}\right)+n$ is \quad (a)~$n^{3/5}$; (b)~$n$; (c)~$n^{5/3}$; (d)~$n^{\log_5 3}$; (e)~$nlogn$.
\end{pastexam}

\begin{pastsoln}
答案是 (b)，$T(n) = \Theta(n)$。

用主定理（資結書定理~2.4）：$a = 3$、$b = 5$、$f(n) = n$。 $n^{\log_b a} = n^{\log_5 3}$，$\log_5 3 \approx 0.683 < 1$。
取 $\varepsilon = 1 - \log_5 3 > 0$，$f(n) = n = \Omega(n^{\log_5 3 + \varepsilon})$；正則條件
$a f(n/b) = 3 \cdot \frac{n}{5} = \frac{3}{5} n \le c f(n)$，取 $c = \frac{3}{5} < 1$ 成立。 屬於情況 3，$T(n) = \Theta(f(n)) = \Theta(n)$。

直觀看遞迴樹：第 $i$ 層的總工作是 $(3/5)^i n$，等比遞減，總和不超過 $\frac{1}{1 - 3/5} n = 2.5n$，由第一層主導。

逐項檢查：(a) $n^{3/5}$ 比 $f(n) = n$ 還小，光是最上層就要 $n$，錯；(b) 對；(c) $n^{5/3}$ 太大，錯；(d) $n^{\log_5 3}$ 是只算葉子
的量（情況 1 的答案），但這裡 $f(n)$ 比它大，錯；(e) $n \lg n$ 是每層工作相同（情況 2）時才有，這裡每層遞減，錯。
\end{pastsoln}
